\clearpage\begingroup\NoSourceHead\SourceLeaf{21}
\begin{center}
{\Large\textls[150]{BRIEFWECHSEL.}}\\[1.5em]
[A.]\\[.8em]
[PRINZIPIEN DER MECHANIK.]\\[1.2em]
[1.]\\[.7em]
\textsc{Gauss} an \textsc{Gerling}.\quad Göttingen, 27. Januar 1834.
\end{center}
\noindent .\enspace.\enspace.\enspace.\enspace.\enspace.\enspace.\enspace.\enspace.\enspace.\enspace.\enspace.\enspace.\enspace.\enspace.\enspace.\enspace.
\par\noindent\hspace*{\parindent}Meine Ansicht ist stets gewesen:
\begin{list}{}{\setlength{\leftmargin}{1.5em}\setlength{\labelwidth}{1.2em}\setlength{\labelsep}{.3em}\setlength{\itemsep}{0pt}\setlength{\parsep}{0pt}\setlength{\topsep}{0pt}}
\item[1)] Dass durch das Ausgehen von \textsc{d’Alemberts} Grundprincip alle Dunkelheit
vollkommen wegfällt.
\item[2)] Dass ohne dieses Grundprincip genügende Klarheit überall nicht zu erhalten
steht.
\end{list}
\noindent Das Grundprincip selbst lässt sich aber wohl treffender einkleiden, als in den
mechanischen Lehrbüchern gewöhnlich geschieht. Ich würde es etwa so aussprechen:
»Wenn ein System von Kräften einem andern System von Kräften
in Beziehung auf ein wie immer verknüpftes System von Materiellem \textls[120]{statisch}
aequivalent ist, so ist es ebendemselben, in derselben Beziehung, auch \textls[120]{mechanisch}
aequivalent«.
\par\noindent .\enspace.\enspace.\enspace.\enspace.\enspace.\enspace.\enspace.\enspace.\enspace.\enspace.\enspace.\enspace.\enspace.\enspace.\enspace.\enspace.
\par\clearpage\endgroup
